% =============================================================================
% fig4-6-cfg-regs-fsm.tex --- 图 4.6 配置寄存器与 6 层 FSM (§4.7)
%
% v5.1 修订: 沿 fig3-5 的 counter × 模块矩阵风格,
%   - 上部: cfg_regs 字段分组 (loop bounds / buffer addrs / quant / schedule)
%   - 主体: counter × 模块矩阵 (同 fig3-5)
%   - 下部: DFE / desc_fifo / sequencer 支撑层 (CFG_WRITE 通路)
% =============================================================================
\documentclass[border=4pt]{standalone}
% =============================================================================
% _tikz-style.tex --- FLUX_CNN 论文 TikZ 共享样式包
%
% 用法（每张 figX-Y.tex 头部 \input 本文件）:
%   \documentclass[border=4pt,varwidth=18cm]{standalone}
%   % =============================================================================
% _tikz-style.tex --- FLUX_CNN 论文 TikZ 共享样式包
%
% 用法（每张 figX-Y.tex 头部 \input 本文件）:
%   \documentclass[border=4pt,varwidth=18cm]{standalone}
%   % =============================================================================
% _tikz-style.tex --- FLUX_CNN 论文 TikZ 共享样式包
%
% 用法（每张 figX-Y.tex 头部 \input 本文件）:
%   \documentclass[border=4pt,varwidth=18cm]{standalone}
%   \input{_tikz-style.tex}
%   \begin{document}
%   \begin{tikzpicture}[<额外局部样式>]
%     ...
%   \end{tikzpicture}
%   \end{document}
%
% 风格基线: IEEE 期刊配色 / 白底黑字 / Times New Roman + Microsoft YaHei /
%           禁卡通 / 禁渐变 / 禁阴影 / 禁 3D / 禁 emoji
% 与 figures-rendered/_style.py (matplotlib) 配色一致, 保证全文图风格统一.
% =============================================================================

% ---- 字体（XeLaTeX 必须）----
% 中文字体：Microsoft YaHei（保持原配置，本次只调整英文）.
% 英文字体：Times New Roman, 把 \sffamily / \rmfamily 同时映射到 Times,
%           保证 TikZ 节点无论用哪个 family, 拉丁字符都是 Times New Roman.
% 数学公式（mathmode）保留 latin modern math, 与 Times 视觉差异可接受.
\usepackage{amsmath}
\usepackage{amssymb}
\usepackage{xeCJK}
\setCJKmainfont{Microsoft YaHei}
\setCJKsansfont{Microsoft YaHei}
\setmainfont{Times New Roman}
\setsansfont{Times New Roman}

% ---- 颜色（与 figures-rendered/_style.py PALETTE='default' 同款）----
\usepackage{xcolor}
\definecolor{coBaseline}{HTML}{1f4e79}  % 深蓝 - 主体强调 / 基线
\definecolor{coImproved}{HTML}{2e7d32}  % 深绿 - 提升 / 高亮
\definecolor{coThird}   {HTML}{b8860b}  % 深金 - 第三系列
\definecolor{coFourth}  {HTML}{6a4c93}  % 深紫 - 第四系列
\definecolor{coAnno}    {HTML}{404040}  % 学术深灰 - 注记
\definecolor{coGuide}   {HTML}{808080}  % 辅助线
\definecolor{coGrid}    {HTML}{cccccc}  % 网格
\definecolor{coAccent}  {HTML}{c00000}  % 强调红 - 仅用于关键警示 / 跨节点
% 浅色填充（CMYK 转 RGB 近似）
\definecolor{flBlue}    {RGB}{226,238,247}  % 淡蓝 CMYK 30/10/0/0 ≈
\definecolor{flYellow}  {RGB}{252,242,219}  % 淡黄 CMYK 0/10/30/0 ≈
\definecolor{flGray}    {RGB}{235,235,235}  % 淡灰 CMYK 0/0/0/10 ≈
\definecolor{flGreen}   {RGB}{226,239,224}  % 浅绿 - 高亮区域

% ---- TikZ + 常用库 ----
\usepackage{tikz}
\usetikzlibrary{
    arrows.meta,
    positioning,
    calc,
    fit,
    backgrounds,
    decorations.pathreplacing,
    decorations.pathmorphing,
    shapes.geometric,
    shapes.misc,
    shapes.multipart,
    matrix,
    chains,
    shadows.blur,    % 不用阴影 (禁用), 但 fit 节点偶尔用 inner sep
    patterns
}

% ---- 全局 TikZ 默认 ----
\tikzset{
    every picture/.style={
        line cap=round,
        line join=round,
        >={Stealth[length=2.4mm, width=1.6mm]},
        font=\sffamily\small,
    },
    % ---- 节点基础类型 ----
    %
    % 主体方框 (深色边框, 白底) - 用于一般模块
    main box/.style={
        draw=coAnno,
        line width=0.7pt,
        rectangle,
        rounded corners=1.5pt,
        minimum width=20mm,
        minimum height=8mm,
        align=center,
        inner sep=2pt,
    },
    % 强调方框 (深蓝填充 alpha)
    accent box/.style={
        main box,
        fill=flBlue,
    },
    % 计算单元 / PE (淡黄填充)
    pe box/.style={
        main box,
        fill=flYellow,
    },
    % SRAM / 存储 (淡灰填充)
    mem box/.style={
        main box,
        fill=flGray,
    },
    % 控制单元 / FSM (浅绿填充)
    ctrl box/.style={
        main box,
        fill=flGreen,
    },
    % 大容器 (圆角矩形, 虚线边框, 用于 multicore_top / NUMA 节点等)
    container box/.style={
        draw=coAnno,
        line width=0.7pt,
        rounded corners=2pt,
        rectangle,
        dashed,
        inner sep=3mm,
    },
    container solid/.style={
        container box,
        solid,
    },
    %
    % ---- 箭头类型 ----
    %
    flow/.style={
        -{Stealth[length=2.4mm,width=1.6mm]},
        line width=0.7pt,
        draw=coAnno,
    },
    flow thick/.style={
        flow,
        line width=1.2pt,
    },
    flow dashed/.style={
        flow,
        dashed,
    },
    flow accent/.style={
        -{Stealth[length=2.4mm,width=1.6mm]},
        line width=1.0pt,
        draw=coAccent,
        dashed,
    },
    flow data/.style={
        -{Stealth[length=2.4mm,width=1.6mm]},
        line width=0.9pt,
        draw=coBaseline,
    },
    flow weight/.style={
        -{Stealth[length=2.4mm,width=1.6mm]},
        line width=0.9pt,
        draw=coImproved,
    },
    flow psum/.style={
        -{Stealth[length=2.4mm,width=1.6mm]},
        line width=0.9pt,
        draw=coThird,
    },
    bus/.style={
        -{Stealth[length=3mm,width=2mm]},
        line width=1.6pt,
        draw=coAnno,
    },
    handshake/.style={
        <->,
        line width=0.7pt,
        draw=coAnno,
    },
    %
    % ---- 标注 / 文字 ----
    %
    label small/.style={
        font=\sffamily\footnotesize,
        text=coAnno,
        align=center,
    },
    label tiny/.style={
        font=\sffamily\scriptsize,
        text=coAnno,
        align=center,
    },
    label bold/.style={
        font=\sffamily\bfseries\small,
        text=black,
        align=center,
    },
    formula/.style={
        font=\small,
        align=center,
    },
    % 边线上的标注 (背景白底, 不挡线)
    edge label/.style={
        font=\sffamily\scriptsize,
        text=coAnno,
        fill=white,
        inner sep=1.2pt,
        align=center,
    },
    edge label accent/.style={
        edge label,
        text=coAccent,
        font=\sffamily\bfseries\scriptsize,
    },
    %
    % ---- 图例 ----
    %
    legend box/.style={
        draw=coAnno,
        line width=0.4pt,
        fill=white,
        rounded corners=1pt,
        inner sep=2.5pt,
        font=\sffamily\scriptsize,
        text=coAnno,
        align=left,
    },
    %
    % ---- 网格 / 阵列 (用于 PE 阵列示意) ----
    %
    grid cell/.style={
        draw=coAnno,
        line width=0.3pt,
        minimum width=4mm,
        minimum height=4mm,
        inner sep=0pt,
        rectangle,
    },
    grid cell light/.style={
        grid cell,
        fill=flGray,
    },
    grid cell accent/.style={
        grid cell,
        fill=flBlue,
    },
    %
    % ---- 时序波形 (用于内嵌时序示意时, 非主要; 主要用 WaveDrom) ----
    %
    wave hi/.style={line width=1pt, draw=coBaseline},
    wave lo/.style={line width=1pt, draw=coBaseline},
}

% ---- 通用辅助宏 ----

% 章节图标题命令 (备选, 一般由论文 caption 承担, 不在 TikZ 内画)
% \newcommand{\figtitle}[1]{\node[font=\bfseries\small,text=black] {#1};}

% 端口标记: 在节点边缘画个小圆点表示端口
\newcommand{\port}[2][black]{\fill[#1] (#2) circle (0.5mm);}

% 中文/英文双行标签: \biLabel{中文}{English}
\newcommand{\biLabel}[2]{\shortstack{\sffamily #1\\\itshape\footnotesize #2}}

% =============================================================================
% 使用约定 (写 figX-Y.tex 时遵守):
%
% 1. 不在 TikZ 内画"图 X.Y"标题——由 Word caption 承担
% 2. 节点用上述预定义样式 (main box / accent box / pe box / mem box / ctrl box / container box)
%    避免每张图重复定义节点样式
% 3. 数据流统一用 flow data (蓝) / flow weight (绿) / flow psum (金) 三色编码
%    握手 / 控制流用 flow / flow dashed (灰)
%    强调 / 警示用 flow accent (红虚线), 慎用
% 4. 文字标注用 label small / label tiny / label bold / edge label
% 5. 图例统一放右下角, 用 legend box 样式
% 6. 中英混排时: 中文用 \sffamily (思源黑体), 英文数学用 \rmfamily (Times New Roman)
% =============================================================================

%   \begin{document}
%   \begin{tikzpicture}[<额外局部样式>]
%     ...
%   \end{tikzpicture}
%   \end{document}
%
% 风格基线: IEEE 期刊配色 / 白底黑字 / Times New Roman + Microsoft YaHei /
%           禁卡通 / 禁渐变 / 禁阴影 / 禁 3D / 禁 emoji
% 与 figures-rendered/_style.py (matplotlib) 配色一致, 保证全文图风格统一.
% =============================================================================

% ---- 字体（XeLaTeX 必须）----
% 中文字体：Microsoft YaHei（保持原配置，本次只调整英文）.
% 英文字体：Times New Roman, 把 \sffamily / \rmfamily 同时映射到 Times,
%           保证 TikZ 节点无论用哪个 family, 拉丁字符都是 Times New Roman.
% 数学公式（mathmode）保留 latin modern math, 与 Times 视觉差异可接受.
\usepackage{amsmath}
\usepackage{amssymb}
\usepackage{xeCJK}
\setCJKmainfont{Microsoft YaHei}
\setCJKsansfont{Microsoft YaHei}
\setmainfont{Times New Roman}
\setsansfont{Times New Roman}

% ---- 颜色（与 figures-rendered/_style.py PALETTE='default' 同款）----
\usepackage{xcolor}
\definecolor{coBaseline}{HTML}{1f4e79}  % 深蓝 - 主体强调 / 基线
\definecolor{coImproved}{HTML}{2e7d32}  % 深绿 - 提升 / 高亮
\definecolor{coThird}   {HTML}{b8860b}  % 深金 - 第三系列
\definecolor{coFourth}  {HTML}{6a4c93}  % 深紫 - 第四系列
\definecolor{coAnno}    {HTML}{404040}  % 学术深灰 - 注记
\definecolor{coGuide}   {HTML}{808080}  % 辅助线
\definecolor{coGrid}    {HTML}{cccccc}  % 网格
\definecolor{coAccent}  {HTML}{c00000}  % 强调红 - 仅用于关键警示 / 跨节点
% 浅色填充（CMYK 转 RGB 近似）
\definecolor{flBlue}    {RGB}{226,238,247}  % 淡蓝 CMYK 30/10/0/0 ≈
\definecolor{flYellow}  {RGB}{252,242,219}  % 淡黄 CMYK 0/10/30/0 ≈
\definecolor{flGray}    {RGB}{235,235,235}  % 淡灰 CMYK 0/0/0/10 ≈
\definecolor{flGreen}   {RGB}{226,239,224}  % 浅绿 - 高亮区域

% ---- TikZ + 常用库 ----
\usepackage{tikz}
\usetikzlibrary{
    arrows.meta,
    positioning,
    calc,
    fit,
    backgrounds,
    decorations.pathreplacing,
    decorations.pathmorphing,
    shapes.geometric,
    shapes.misc,
    shapes.multipart,
    matrix,
    chains,
    shadows.blur,    % 不用阴影 (禁用), 但 fit 节点偶尔用 inner sep
    patterns
}

% ---- 全局 TikZ 默认 ----
\tikzset{
    every picture/.style={
        line cap=round,
        line join=round,
        >={Stealth[length=2.4mm, width=1.6mm]},
        font=\sffamily\small,
    },
    % ---- 节点基础类型 ----
    %
    % 主体方框 (深色边框, 白底) - 用于一般模块
    main box/.style={
        draw=coAnno,
        line width=0.7pt,
        rectangle,
        rounded corners=1.5pt,
        minimum width=20mm,
        minimum height=8mm,
        align=center,
        inner sep=2pt,
    },
    % 强调方框 (深蓝填充 alpha)
    accent box/.style={
        main box,
        fill=flBlue,
    },
    % 计算单元 / PE (淡黄填充)
    pe box/.style={
        main box,
        fill=flYellow,
    },
    % SRAM / 存储 (淡灰填充)
    mem box/.style={
        main box,
        fill=flGray,
    },
    % 控制单元 / FSM (浅绿填充)
    ctrl box/.style={
        main box,
        fill=flGreen,
    },
    % 大容器 (圆角矩形, 虚线边框, 用于 multicore_top / NUMA 节点等)
    container box/.style={
        draw=coAnno,
        line width=0.7pt,
        rounded corners=2pt,
        rectangle,
        dashed,
        inner sep=3mm,
    },
    container solid/.style={
        container box,
        solid,
    },
    %
    % ---- 箭头类型 ----
    %
    flow/.style={
        -{Stealth[length=2.4mm,width=1.6mm]},
        line width=0.7pt,
        draw=coAnno,
    },
    flow thick/.style={
        flow,
        line width=1.2pt,
    },
    flow dashed/.style={
        flow,
        dashed,
    },
    flow accent/.style={
        -{Stealth[length=2.4mm,width=1.6mm]},
        line width=1.0pt,
        draw=coAccent,
        dashed,
    },
    flow data/.style={
        -{Stealth[length=2.4mm,width=1.6mm]},
        line width=0.9pt,
        draw=coBaseline,
    },
    flow weight/.style={
        -{Stealth[length=2.4mm,width=1.6mm]},
        line width=0.9pt,
        draw=coImproved,
    },
    flow psum/.style={
        -{Stealth[length=2.4mm,width=1.6mm]},
        line width=0.9pt,
        draw=coThird,
    },
    bus/.style={
        -{Stealth[length=3mm,width=2mm]},
        line width=1.6pt,
        draw=coAnno,
    },
    handshake/.style={
        <->,
        line width=0.7pt,
        draw=coAnno,
    },
    %
    % ---- 标注 / 文字 ----
    %
    label small/.style={
        font=\sffamily\footnotesize,
        text=coAnno,
        align=center,
    },
    label tiny/.style={
        font=\sffamily\scriptsize,
        text=coAnno,
        align=center,
    },
    label bold/.style={
        font=\sffamily\bfseries\small,
        text=black,
        align=center,
    },
    formula/.style={
        font=\small,
        align=center,
    },
    % 边线上的标注 (背景白底, 不挡线)
    edge label/.style={
        font=\sffamily\scriptsize,
        text=coAnno,
        fill=white,
        inner sep=1.2pt,
        align=center,
    },
    edge label accent/.style={
        edge label,
        text=coAccent,
        font=\sffamily\bfseries\scriptsize,
    },
    %
    % ---- 图例 ----
    %
    legend box/.style={
        draw=coAnno,
        line width=0.4pt,
        fill=white,
        rounded corners=1pt,
        inner sep=2.5pt,
        font=\sffamily\scriptsize,
        text=coAnno,
        align=left,
    },
    %
    % ---- 网格 / 阵列 (用于 PE 阵列示意) ----
    %
    grid cell/.style={
        draw=coAnno,
        line width=0.3pt,
        minimum width=4mm,
        minimum height=4mm,
        inner sep=0pt,
        rectangle,
    },
    grid cell light/.style={
        grid cell,
        fill=flGray,
    },
    grid cell accent/.style={
        grid cell,
        fill=flBlue,
    },
    %
    % ---- 时序波形 (用于内嵌时序示意时, 非主要; 主要用 WaveDrom) ----
    %
    wave hi/.style={line width=1pt, draw=coBaseline},
    wave lo/.style={line width=1pt, draw=coBaseline},
}

% ---- 通用辅助宏 ----

% 章节图标题命令 (备选, 一般由论文 caption 承担, 不在 TikZ 内画)
% \newcommand{\figtitle}[1]{\node[font=\bfseries\small,text=black] {#1};}

% 端口标记: 在节点边缘画个小圆点表示端口
\newcommand{\port}[2][black]{\fill[#1] (#2) circle (0.5mm);}

% 中文/英文双行标签: \biLabel{中文}{English}
\newcommand{\biLabel}[2]{\shortstack{\sffamily #1\\\itshape\footnotesize #2}}

% =============================================================================
% 使用约定 (写 figX-Y.tex 时遵守):
%
% 1. 不在 TikZ 内画"图 X.Y"标题——由 Word caption 承担
% 2. 节点用上述预定义样式 (main box / accent box / pe box / mem box / ctrl box / container box)
%    避免每张图重复定义节点样式
% 3. 数据流统一用 flow data (蓝) / flow weight (绿) / flow psum (金) 三色编码
%    握手 / 控制流用 flow / flow dashed (灰)
%    强调 / 警示用 flow accent (红虚线), 慎用
% 4. 文字标注用 label small / label tiny / label bold / edge label
% 5. 图例统一放右下角, 用 legend box 样式
% 6. 中英混排时: 中文用 \sffamily (思源黑体), 英文数学用 \rmfamily (Times New Roman)
% =============================================================================

%   \begin{document}
%   \begin{tikzpicture}[<额外局部样式>]
%     ...
%   \end{tikzpicture}
%   \end{document}
%
% 风格基线: IEEE 期刊配色 / 白底黑字 / Times New Roman + Microsoft YaHei /
%           禁卡通 / 禁渐变 / 禁阴影 / 禁 3D / 禁 emoji
% 与 figures-rendered/_style.py (matplotlib) 配色一致, 保证全文图风格统一.
% =============================================================================

% ---- 字体（XeLaTeX 必须）----
% 中文字体：Microsoft YaHei（保持原配置，本次只调整英文）.
% 英文字体：Times New Roman, 把 \sffamily / \rmfamily 同时映射到 Times,
%           保证 TikZ 节点无论用哪个 family, 拉丁字符都是 Times New Roman.
% 数学公式（mathmode）保留 latin modern math, 与 Times 视觉差异可接受.
\usepackage{amsmath}
\usepackage{amssymb}
\usepackage{xeCJK}
\setCJKmainfont{Microsoft YaHei}
\setCJKsansfont{Microsoft YaHei}
\setmainfont{Times New Roman}
\setsansfont{Times New Roman}

% ---- 颜色（与 figures-rendered/_style.py PALETTE='default' 同款）----
\usepackage{xcolor}
\definecolor{coBaseline}{HTML}{1f4e79}  % 深蓝 - 主体强调 / 基线
\definecolor{coImproved}{HTML}{2e7d32}  % 深绿 - 提升 / 高亮
\definecolor{coThird}   {HTML}{b8860b}  % 深金 - 第三系列
\definecolor{coFourth}  {HTML}{6a4c93}  % 深紫 - 第四系列
\definecolor{coAnno}    {HTML}{404040}  % 学术深灰 - 注记
\definecolor{coGuide}   {HTML}{808080}  % 辅助线
\definecolor{coGrid}    {HTML}{cccccc}  % 网格
\definecolor{coAccent}  {HTML}{c00000}  % 强调红 - 仅用于关键警示 / 跨节点
% 浅色填充（CMYK 转 RGB 近似）
\definecolor{flBlue}    {RGB}{226,238,247}  % 淡蓝 CMYK 30/10/0/0 ≈
\definecolor{flYellow}  {RGB}{252,242,219}  % 淡黄 CMYK 0/10/30/0 ≈
\definecolor{flGray}    {RGB}{235,235,235}  % 淡灰 CMYK 0/0/0/10 ≈
\definecolor{flGreen}   {RGB}{226,239,224}  % 浅绿 - 高亮区域

% ---- TikZ + 常用库 ----
\usepackage{tikz}
\usetikzlibrary{
    arrows.meta,
    positioning,
    calc,
    fit,
    backgrounds,
    decorations.pathreplacing,
    decorations.pathmorphing,
    shapes.geometric,
    shapes.misc,
    shapes.multipart,
    matrix,
    chains,
    shadows.blur,    % 不用阴影 (禁用), 但 fit 节点偶尔用 inner sep
    patterns
}

% ---- 全局 TikZ 默认 ----
\tikzset{
    every picture/.style={
        line cap=round,
        line join=round,
        >={Stealth[length=2.4mm, width=1.6mm]},
        font=\sffamily\small,
    },
    % ---- 节点基础类型 ----
    %
    % 主体方框 (深色边框, 白底) - 用于一般模块
    main box/.style={
        draw=coAnno,
        line width=0.7pt,
        rectangle,
        rounded corners=1.5pt,
        minimum width=20mm,
        minimum height=8mm,
        align=center,
        inner sep=2pt,
    },
    % 强调方框 (深蓝填充 alpha)
    accent box/.style={
        main box,
        fill=flBlue,
    },
    % 计算单元 / PE (淡黄填充)
    pe box/.style={
        main box,
        fill=flYellow,
    },
    % SRAM / 存储 (淡灰填充)
    mem box/.style={
        main box,
        fill=flGray,
    },
    % 控制单元 / FSM (浅绿填充)
    ctrl box/.style={
        main box,
        fill=flGreen,
    },
    % 大容器 (圆角矩形, 虚线边框, 用于 multicore_top / NUMA 节点等)
    container box/.style={
        draw=coAnno,
        line width=0.7pt,
        rounded corners=2pt,
        rectangle,
        dashed,
        inner sep=3mm,
    },
    container solid/.style={
        container box,
        solid,
    },
    %
    % ---- 箭头类型 ----
    %
    flow/.style={
        -{Stealth[length=2.4mm,width=1.6mm]},
        line width=0.7pt,
        draw=coAnno,
    },
    flow thick/.style={
        flow,
        line width=1.2pt,
    },
    flow dashed/.style={
        flow,
        dashed,
    },
    flow accent/.style={
        -{Stealth[length=2.4mm,width=1.6mm]},
        line width=1.0pt,
        draw=coAccent,
        dashed,
    },
    flow data/.style={
        -{Stealth[length=2.4mm,width=1.6mm]},
        line width=0.9pt,
        draw=coBaseline,
    },
    flow weight/.style={
        -{Stealth[length=2.4mm,width=1.6mm]},
        line width=0.9pt,
        draw=coImproved,
    },
    flow psum/.style={
        -{Stealth[length=2.4mm,width=1.6mm]},
        line width=0.9pt,
        draw=coThird,
    },
    bus/.style={
        -{Stealth[length=3mm,width=2mm]},
        line width=1.6pt,
        draw=coAnno,
    },
    handshake/.style={
        <->,
        line width=0.7pt,
        draw=coAnno,
    },
    %
    % ---- 标注 / 文字 ----
    %
    label small/.style={
        font=\sffamily\footnotesize,
        text=coAnno,
        align=center,
    },
    label tiny/.style={
        font=\sffamily\scriptsize,
        text=coAnno,
        align=center,
    },
    label bold/.style={
        font=\sffamily\bfseries\small,
        text=black,
        align=center,
    },
    formula/.style={
        font=\small,
        align=center,
    },
    % 边线上的标注 (背景白底, 不挡线)
    edge label/.style={
        font=\sffamily\scriptsize,
        text=coAnno,
        fill=white,
        inner sep=1.2pt,
        align=center,
    },
    edge label accent/.style={
        edge label,
        text=coAccent,
        font=\sffamily\bfseries\scriptsize,
    },
    %
    % ---- 图例 ----
    %
    legend box/.style={
        draw=coAnno,
        line width=0.4pt,
        fill=white,
        rounded corners=1pt,
        inner sep=2.5pt,
        font=\sffamily\scriptsize,
        text=coAnno,
        align=left,
    },
    %
    % ---- 网格 / 阵列 (用于 PE 阵列示意) ----
    %
    grid cell/.style={
        draw=coAnno,
        line width=0.3pt,
        minimum width=4mm,
        minimum height=4mm,
        inner sep=0pt,
        rectangle,
    },
    grid cell light/.style={
        grid cell,
        fill=flGray,
    },
    grid cell accent/.style={
        grid cell,
        fill=flBlue,
    },
    %
    % ---- 时序波形 (用于内嵌时序示意时, 非主要; 主要用 WaveDrom) ----
    %
    wave hi/.style={line width=1pt, draw=coBaseline},
    wave lo/.style={line width=1pt, draw=coBaseline},
}

% ---- 通用辅助宏 ----

% 章节图标题命令 (备选, 一般由论文 caption 承担, 不在 TikZ 内画)
% \newcommand{\figtitle}[1]{\node[font=\bfseries\small,text=black] {#1};}

% 端口标记: 在节点边缘画个小圆点表示端口
\newcommand{\port}[2][black]{\fill[#1] (#2) circle (0.5mm);}

% 中文/英文双行标签: \biLabel{中文}{English}
\newcommand{\biLabel}[2]{\shortstack{\sffamily #1\\\itshape\footnotesize #2}}

% =============================================================================
% 使用约定 (写 figX-Y.tex 时遵守):
%
% 1. 不在 TikZ 内画"图 X.Y"标题——由 Word caption 承担
% 2. 节点用上述预定义样式 (main box / accent box / pe box / mem box / ctrl box / container box)
%    避免每张图重复定义节点样式
% 3. 数据流统一用 flow data (蓝) / flow weight (绿) / flow psum (金) 三色编码
%    握手 / 控制流用 flow / flow dashed (灰)
%    强调 / 警示用 flow accent (红虚线), 慎用
% 4. 文字标注用 label small / label tiny / label bold / edge label
% 5. 图例统一放右下角, 用 legend box 样式
% 6. 中英混排时: 中文用 \sffamily (思源黑体), 英文数学用 \rmfamily (Times New Roman)
% =============================================================================


\begin{document}
\begin{tikzpicture}[
    every node/.style={font=\sffamily\footnotesize},
    cfg group/.style={
        draw=coAnno, line width=0.5pt, rounded corners=1pt,
        fill=flYellow, minimum width=34mm, minimum height=10mm,
        align=center, inner sep=2pt,
        font=\sffamily\scriptsize\bfseries,
    },
    cfg main/.style={
        draw=coAnno, line width=0.8pt, rounded corners=2pt,
        fill=flGreen, minimum width=28mm, minimum height=10mm,
        align=center, inner sep=2pt,
        font=\sffamily\small\bfseries,
    },
    cell/.style={
        draw=coAnno, line width=0.4pt, rectangle,
        minimum width=22mm, minimum height=8mm,
        inner sep=0pt, align=center,
        font=\sffamily\scriptsize,
    },
    cell head/.style={cell, fill=flGreen, font=\sffamily\scriptsize\bfseries,},
    cell row/.style={cell, fill=flYellow, font=\sffamily\scriptsize\bfseries,},
    cell main/.style={cell, fill=flBlue, font=\sffamily\scriptsize,},
    cell sub/.style={cell, fill=white, font=\sffamily\scriptsize, text=coGuide,},
    support/.style={
        draw=coAnno, line width=0.6pt, rounded corners=1pt,
        fill=flGray, minimum width=24mm, minimum height=8mm,
        align=center, inner sep=2pt,
        font=\sffamily\scriptsize\bfseries,
    },
    flow ctrl/.style={
        -{Stealth[length=2.5mm,width=1.8mm]},
        line width=0.7pt, draw=coAnno,
    },
    flow cfgw/.style={
        -{Stealth[length=2.8mm,width=2mm]},
        line width=1.0pt, draw=coAccent,
    },
    section title/.style={
        font=\sffamily\small\bfseries, text=coAnno, anchor=west,
    },
]

% =============================================================================
% Section 1: cfg_regs 字段分组 (左上)
% =============================================================================
\node[section title] at (-0.5, 6.3)
    {cfg\_regs 字段分组 \ (host AXI-Lite + DFE CFG\_WRITE 双写口)};

\node[cfg group] (g1) at (1.5, 5.3) {Loop bounds\\\scriptsize $H_{out}/W_{out}/C_{in}/C_{out}$\\\scriptsize $K/$stride$/$pad};
\node[cfg group] (g2) at (5.5, 5.3) {Buffer addrs\\\scriptsize IFB / WB / OFB\\\scriptsize bias / shortcut};
\node[cfg group] (g3) at (9.5, 5.3) {Quant params\\\scriptsize shift\_amt\\\scriptsize sat\_min/max};
\node[cfg group] (g4) at (13.0, 5.3) {Schedule\\\scriptsize $n_{cs}/n_{tile}/n_{cins}$\\\scriptsize tile width};

\node[cfg main] (cfg) at (7.0, 3.6) {cfg\_regs};

\foreach \g/\dx in {g1/-3.6, g2/-0.6, g3/2.6, g4/5.6}{
    \draw[flow ctrl] (\g.south) -- ($(cfg.north)+(\dx*0.18,0)$);
}

% =============================================================================
% Section 2: counter × 模块矩阵 (主体, 与 fig3-5 同风格)
% =============================================================================
\node[section title] at (-0.5, 2.4)
    {6 级 counter $\times$ 5 模块矩阵 \ (本地 counter 推进, valid-ready 握手汇流)};

\def\colw{2.2}
\def\rowh{0.8}
\def\xz{1.0}
\def\yz{1.4}

\foreach \c/\mod in {0/line\_buffer, 1/mac\_array, 2/parf\_accum, 3/ofb\_writer, 4/wgt\_buffer}{
    \node[cell head] at ({\xz+\c*\colw}, \yz) {\mod};
}

% Row 0: cs
\node[cell row]  at ({\xz-\colw}, {\yz-\rowh}) {cs\\\scriptsize cout slice};
\node[cell sub]  at ({\xz+0*\colw}, {\yz-\rowh}) {-};
\node[cell main] at ({\xz+1*\colw}, {\yz-\rowh}) {WRF 切片};
\node[cell main] at ({\xz+2*\colw}, {\yz-\rowh}) {累加结束};
\node[cell main] at ({\xz+3*\colw}, {\yz-\rowh}) {bias 切换};
\node[cell main] at ({\xz+4*\colw}, {\yz-\rowh}) {WB load};

% Row 1: yout
\node[cell row]  at ({\xz-\colw}, {\yz-2*\rowh}) {yout\\\scriptsize 输出行};
\node[cell main] at ({\xz+0*\colw}, {\yz-2*\rowh}) {ring 行步进};
\node[cell sub]  at ({\xz+1*\colw}, {\yz-2*\rowh}) {-};
\node[cell sub]  at ({\xz+2*\colw}, {\yz-2*\rowh}) {-};
\node[cell main] at ({\xz+3*\colw}, {\yz-2*\rowh}) {OFB 行步进};
\node[cell sub]  at ({\xz+4*\colw}, {\yz-2*\rowh}) {-};

% Row 2: tile
\node[cell row]  at ({\xz-\colw}, {\yz-3*\rowh}) {tile\\\scriptsize $W$ 维 tile};
\node[cell main] at ({\xz+0*\colw}, {\yz-3*\rowh}) {tile addr};
\node[cell sub]  at ({\xz+1*\colw}, {\yz-3*\rowh}) {-};
\node[cell main] at ({\xz+2*\colw}, {\yz-3*\rowh}) {cs 维发射};
\node[cell main] at ({\xz+3*\colw}, {\yz-3*\rowh}) {写 OFM};
\node[cell sub]  at ({\xz+4*\colw}, {\yz-3*\rowh}) {-};

% Row 3: cins
\node[cell row]  at ({\xz-\colw}, {\yz-4*\rowh}) {cins\\\scriptsize cin slice};
\node[cell main] at ({\xz+0*\colw}, {\yz-4*\rowh}) {cin 步进};
\node[cell main] at ({\xz+1*\colw}, {\yz-4*\rowh}) {ARF 切换};
\node[cell main] at ({\xz+2*\colw}, {\yz-4*\rowh}) {psum 累加};
\node[cell sub]  at ({\xz+3*\colw}, {\yz-4*\rowh}) {-};
\node[cell main] at ({\xz+4*\colw}, {\yz-4*\rowh}) {WRF 重载};

% Row 4: round
\node[cell row]  at ({\xz-\colw}, {\yz-5*\rowh}) {round\\\scriptsize $K^2$ 轮次};
\node[cell main] at ({\xz+0*\colw}, {\yz-5*\rowh}) {ARF 读};
\node[cell main] at ({\xz+1*\colw}, {\yz-5*\rowh}) {PE 取权重};
\node[cell sub]  at ({\xz+2*\colw}, {\yz-5*\rowh}) {-};
\node[cell sub]  at ({\xz+3*\colw}, {\yz-5*\rowh}) {-};
\node[cell main] at ({\xz+4*\colw}, {\yz-5*\rowh}) {WRF 读};

% Row 5: pos
\node[cell row]  at ({\xz-\colw}, {\yz-6*\rowh}) {pos\\\scriptsize PE 阵列位};
\node[cell main] at ({\xz+0*\colw}, {\yz-6*\rowh}) {ARF 读 1 拍};
\node[cell main] at ({\xz+1*\colw}, {\yz-6*\rowh}) {PE 发射 1 拍};
\node[cell sub]  at ({\xz+2*\colw}, {\yz-6*\rowh}) {-};
\node[cell sub]  at ({\xz+3*\colw}, {\yz-6*\rowh}) {-};
\node[cell sub]  at ({\xz+4*\colw}, {\yz-6*\rowh}) {-};

% 左侧大括号
\draw[decorate, decoration={brace, amplitude=4pt, mirror},
      line width=0.6pt, draw=coAccent]
    ({\xz-\colw-1.25}, {\yz-\rowh+0.4})
    -- ({\xz-\colw-1.25}, {\yz-6*\rowh-0.4})
    node[midway, left=4pt, font=\sffamily\scriptsize\bfseries, text=coAccent,
         align=center]
    {6 级嵌套\\外 $\to$ 内};

% =============================================================================
% Section 3: 底部 DFE / desc_fifo / sequencer 支撑层
% =============================================================================
\node[section title] at (-0.5, -5.0)
    {CFG\_WRITE 通路 \ (DFE → desc\_fifo → sequencer → cfg\_regs seq\_w port)};

\node[support] (dfe)  at (2.0, -6.0) {DFE\\\scriptsize desc\_fetch\_engine};
\node[support] (fifo) at (5.0, -6.0) {desc\_fifo\\\scriptsize 128 entries};
\node[support] (seq)  at (8.0, -6.0) {sequencer};
\node[support] (cfgw) at (11.0, -6.0) {cfg\_regs\\\scriptsize seq\_w port};

\draw[flow ctrl] (dfe.east)  -- (fifo.west);
\draw[flow ctrl] (fifo.east) -- (seq.west);
\draw[flow cfgw] (seq.east)  -- (cfgw.west)
    node[edge label, midway, above=1pt, text=coAccent]{CFG\_WRITE};

% cfg_regs (上) 与 cfg_regs seq_w (下) 用虚线连接表示同一物理寄存器组
\draw[flow ctrl, dashed]
    (cfg.east) to[bend left=20] (cfgw.north);

\end{tikzpicture}
\end{document}
